%%%%%%
%
%%%
%
% $Autor: Wings $
% $Datum: 2021-05-14 $
% $Pfad: GitLab/210419KDDInverseKinematic $
% $Dateiname: Introduction
% $Version: 4620 $
%
% !TeX spellcheck = de_GB
% !TeX program = pdflatex
% !BIB program = biber
% !TeX encoding = utf8
%
%%%

\chapter{Hardware Description}
\section{Introduction to Nvidia Jetson Nano}
The NVIDIA® Jetson NanoTM Developer Kit is a maker, learner, and developer AI computer that brings the power of modern artificial intelligence to a low-power, user-friendly platform. Get started quickly with built-in support for a wide range of popular peripherals, add-ons, and ready-to-use projects.

Jetson Nano is supported by the extensive NVIDIA® JetPackTM SDK and comes with the performance and capabilities required to run modern AI workloads as listed below. JetPack includes the following: \textbf{[Reference: https://cdn-reichelt.de/documents/datenblatt/I100/NVIDIA\_JETSON\_NANO\_ENG\\\_MAN.pdf]}
\begin{enumerate}
	\item Full desktop Linux with NVIDIA drivers
	\item AI and Computer Vision libraries and APIs
	\item Developer tools
	\item Documentation and sample code
\end{enumerate}

\textbf{Processing Power}

Jetson Nano has 472 GFLOPs to tackle contemporary AI algorithms. It can run many neural networks in parallel and interpret numerous high-resolution sensors at the same time, making it suitable for NVRs and intelligent gateways.\\
It consists of a NVIDIA Maxwell™ architecture GPU with 128 NVIDIA CUDA® cores running at 0.5 TFLOPs (FP16) and a Quad-core ARM® Cortex®-A57 MPCore processor. A RAM of 4 GB 64-bit LPDDR4 running 1600MHz (25.6 GB/s) coupled with a storage of 16gb (16 GB eMMC 5.1 Flash) provides adequate multi tasking power. \textbf{[Reference:https://www.nvidia.com/en-us/autonomous-machines/embedded-systems/jetson-nano/product-development/]}\\
\textbf{Connectivity}\\
With the help of an add on chip, WiFi (802.11b/g/n) connection is possible. The connection can also be achieved with the help of a Ethernet cable. Based on the specification of the cable used (10/100/1000 BASE-T) the speed may vary.\\
\textbf{Programming}\\
Nividia provides a variety of Linux based GUI's for ease of programming of the Jetson Nano as follows:\\
\begin{enumerate}
	\item  \textbf{NVIDIA Jetpack SDK:}\\ 
	The Jetson software stack begins with NVIDIA JetPack™ SDK, which provides Jetson Linux, developer tools, and CUDA-X accelerated libraries and other NVIDIA technologies.\\
	
	JetPack enables end-to-end acceleration for your AI applications, with NVIDIA TensorRT and cuDNN for accelerated AI inferencing, CUDA for accelerated general computing, VPI for accelerated computer vision and image processing, Jetson Linux API’s for accelerated multimedia, and libArgus and V4l2 for accelerated camera processing.\\
	
	NVIDIA container runtime is also included in JetPack, enabling cloud-native technologies and workflows at the edge. Transform your experience of developing and deploying software by containerizing your AI applications and managing them at scale with cloud-native technologies.\\
	
	Jetson Linux provides the foundation for your applications with a Linux kernel, bootloader, NVIDIA drivers, flashing utilities, sample filesystem, and toolchains for the Jetson platform. It also includes security features, over-the-air update capabilities and much more.\\
	\textbf{[Reference: https://developer.nvidia.com/embedded/develop/software]}
	\item \textbf{NVIDIA TAO \& Pretrained AI Models:}\\
	NVIDIA TAO shortens the time to value by simplifying the time-consuming components of a deep learning workflow, from data preparation to training to optimization. Start with production-ready pre-trained AI models from the NVIDIA NGCTM collection to accelerate development by 10X. These models have been taught to achieve high accuracy in a variety of fields, including computer vision, conversational AI, and others.\\
	
	Data gathering and annotation is a time-consuming and costly procedure. Simulation can assist you in meeting this data requirement in an effective manner. NVIDIA Omniverse Replicator employs simulation to produce synthetic data that is orders of magnitude quicker and less expensive than acquiring real-world data. You can easily produce varied, huge, and accurate datasets for training AI models using Omniverse Replicator.\textbf{[Reference: https://developer.nvidia.com/\\embedded/develop/software]}
	
	\item \textbf{NVIDIA TritionTM:}\\
	The NVIDIA TritonTM Inference Server makes it easier to deploy AI models at scale. Triton Inference Server is open source and provides a single standardized inference platform capable of supporting multi framework model inference in various deployments such as datacenter, cloud, embedded devices, and virtualized environments. It supports a variety of inference questions using powerful batching and scheduling methods, as well as live model changes.\textbf{[Reference: https://developer.nvidia.com/embedded/develop/software]}
	
	\item \textbf{NVIDIA Riva:}\\
	NVIDIA Riva is a fully accelerated SDK for developing multimodal conversational AI apps employing a deep learning pipeline from start to finish. Pretrained conversational AI models, the NVIDIA TAO Toolkit, and optimized end-to-end capabilities for voice, vision, and natural language processing (NLP) applications are all included in the Riva SDK.\textbf{[Reference: https://developer.nvidia.com/\\embedded/develop/software]}
	
	\item \textbf{NVIDIA DeepStream SDK:}\\
	NVIDIA Riva is a fully accelerated SDK for developing multimodal conversational AI apps employing a deep learning pipeline from start to finish. Pretrained conversational AI models, the NVIDIA TAO Toolkit, and optimized end-to-end capabilities for voice, vision, and natural language processing (NLP) applications are all included in the Riva SDK.\textbf{[Reference: https://developer.nvidia.com/\\embedded/develop/software]}
	
	\item \textbf{NVIDIA Isaac:}\\
	NVIDIA Isaac ROS GEMs are hardware-accelerated packages that enable ROS developers to design high-performance solutions on NVIDIA hardware with more ease. NVIDIA Isaac Sim is a scalable robotics simulation program powered by Omniverse. It comes with Replicator, a tool for creating various synthetic datasets for training perception models. Isaac Sim is also a tool for developing, testing, and managing AI-based robots by providing photorealistic, physically correct virtual settings.\textbf{[Reference: https://developer.nvidia.com/\\embedded/develop/software]}
	
\end{enumerate}
\section{Board layout of the Nvidia Jetson Nano}
\begin{figure}[h]
	\centering
	\includegraphics{Images/Board\_Layout}
	\caption{Board Layout}
\end{figure}
This section highlights some of the Jetson Nano Developer Kit interfaces.
\begin{itemize}
	\item Module
	\begin{itemize}
		\item [J501] Slot for a microSD card.
		\item The passive heatsink supports 10W module power usage at 25° C ambient temperature. If your user case requires additional cooling, you can configure the	module to control a system fan.
	\end{itemize}
	\item Carrier Board
	\begin{itemize}
		\item (DS3) Power LED; lights when the developer kit is powered on.
		\item (J2) SO-DIMM connector for Jetson Nano module.
		\item (J6] HDMI and DP connector stack.
		\item (J13) Camera connector; enables use of CSI cameras. Jetson Nano Developer Kit works with IMX219 camera modules, including Leopard Imaging LI-IMX219-MIPIFF-NANO camera module and Raspberry Pi Camera Module V2.
		\item (J15) 4-pin fan control header. Pulse Width Modulation (PWM) output and tachometer input are supported.
		\item (J18) M.2 Key E connector can be used for wireless networking cards; includes	interfaces for PCIe (x1), USB 2.0, UART, I2S, and I2C. To reach J18 you must detach the Jetson Nano module.
		\item (J25) Power jack for 5V~4A power supply. Accepts a 2.1×5.5×9.5 mm plug with positive polarity.
		\item (J28) Micro-USB 2.0 connector; can be used in either of two ways:\\
		If J48 pins are not connected, you can power the developer kit from a 5V~2A
		Micro-USB power supply.\\
		If J48 pins are connected, operates in Device Mode.
		\item (J32 and J33) are each a stack of two USB 3.0 Type A connectors. Each stack is limited
		to 1A total power delivery. All four are connected to the Jetson Nano module via a
		USB 3.0 hub built into the carrier board.
		\item (J38) The Power over Ethernet (POE) header exposes any DC voltage present on J43 Ethernet jack per IEEE 802.3af.
		\item (J40) 8-pin button header; brings out several system power, reset, and force recovery related signals.
	\end{itemize}
\end{itemize}
\textbf{[Reference: https://developer.nvidia.com/\\embedded/develop/software]}
\section{Specs of Nvidia Jetson Nano}
The specifications of the NVIDIA Jetson Nano is as given in the table below:\\
\begin{table}
	\begin{tabular}{||m{5cm}|m{10cm}||}
		\hline
		\textbf{Item} & \textbf{Specification}\\
		\hline
		GPU & NVIDIA Maxwell architecture with 128 NVIDIA CUDA® cores\\
		\hline
		CPU & Quad-core ARM Cortex-A57 MPCore processor\\
		\hline
		Memory & 4 GB 64-bit LPDDR4, 1600MHz 25.6 GB/s\\
		\hline
		Storage & 16 GB eMMC 5.1\\
		\hline
		Video encode & 250MP/sec
		1x 4K @ 30 (HEVC)
		2x 1080p @ 60 (HEVC)
		4x 1080p @ 30 (HEVC)
		4x 720p @ 60 (HEVC)
		9x 720p @ 30 (HEVC)\\
		\hline
		Video decode & 500MP/sec
		1x 4K @ 60 (HEVC)
		2x 4K @ 30 (HEVC)
		4x 1080p @ 60 (HEVC)
		8x 1080p @ 30 (HEVC)
		9x 720p @ 60 (HEVC)\\
		\hline
		Camera & 12 lanes (3x4 or 4x2) MIPI CSI-2 D-PHY 1.1 (1.5 Gb/s per pair)\\
		\hline
		Connectivity & Gigabit Ethernet, M.2 Key E\\
		\hline
		Display & HDMI 2.0 and eDP 1.4\\
		\hline
		USB & 4x USB 3.0, USB 2.0 Micro-B\\
		\hline
		Others & GPIO, I2C, I2S, SPI, UART\\
		\hline
		Mechanical & 69.6 mm x 45 mm
		260-pin edge connector\\
		\hline
	\end{tabular}
\caption{\label{tab:NVIDIA Jetson Nano Specification}Detailed Specifications of the NVIDIA Jetson Nano\\ \textbf{[Reference: https://developer.nvidia.com/embedded/jetson-nano]}}
\end{table}

